\documentclass{article}
\usepackage[utf8]{inputenc} 
\usepackage[spanish]{babel} 
\usepackage{tikz}
\usepackage{graphicx} 
\usepackage{listings}
\usepackage{amsmath}
\usepackage{forest}
\usetikzlibrary{shadows}
\usepackage{xcolor}
\definecolor{codegreen}{rgb}{0,0.6,0}
\definecolor{codegray}{rgb}{0.5,0.5,0.5}
\definecolor{codepurple}{rgb}{0.58,0,0.82}
\definecolor{backcolour}{rgb}{0.95,0.95,0.92}
\lstset{
    backgroundcolor=\color{backcolour},   
    commentstyle=\color{codegreen},
    keywordstyle=\color{magenta},
    numberstyle=\tiny\color{codegray},
    stringstyle=\color{codepurple},
    basicstyle=\ttfamily\footnotesize,
    breaklines=true,                 
    captionpos=b,                    
    keepspaces=true,                 
    numbers=left,                    
    showspaces=false,                
    language=C++
}
\begin{document}
\begin{center}
    \huge \textbf{Dynamic Programming}
    \vspace{0.5cm}
\end{center}
\section*{Introduccion}
Para introducirnos en programacion dinamica podemos tomar en cuenta el ejemplo mas conocido, el cual es la sucesion de fibonacci, para ello tenemos la solucion de forma recursiva:
\begin{equation}
    f(n)=
    \begin{cases}
        0,& si: n=0 \\
        1,& si: n=1 \\
        f(n-1)+f(n-2),& si: n>1
    \end{cases}
\end{equation}
Pero esta solucion es muy poco optima para casos grandes ya que su complejidad es exponencial, es decir $O(2^n)$, por lo cual revisaremos cual es su mayor problema:\par
\begin{forest}
  for tree={
    circle,
    draw=blue!60,
    fill=blue!10,
    thick,
    minimum size=0.5cm,
    inner sep=1pt,
    font=\sffamily\bfseries,
    edge={-stealth, gray!60, line width=1pt},
    l sep=0.5cm, 
    s sep=0.8cm 
  }
  [f(5)
    [f(4)
      [f(3),fill=orange!20 
        [f(2)
        [f(1)]
        [f(0)]
        ]
        [f(1)]
      ]
      [f(2)
        [f(1)]
        [f(0)]
      ]
    ]
    [f(3),fill=orange!20 
      [f(2)
        [f(1)]
        [f(0)]
      ]
      [f(1)]
    ]
  ]
\end{forest}
\par
Como se puede ver el mayor problema de la recursion es que los calculos se repiten, podemos ver que $f(3)$ se calcula dos veces, para numeros mas grandes estos calculos repetidos se multiplican, para ello podemos utilizar un vector el cual guarde los datos ya calculados:
\begin{lstlisting}
#define ll long long
vector<ll>memo(N+1,-1);
ll fib(ll n){
    if(memo[n]!=-1)return memo[n];
    if(n==0||n==1)return n;
    return memo[n]=fib1(n-1)+fib1(n-2);
}
\end{lstlisting}
Utilizando la memorizacion podemo reducir la complejidad exponencail en lineal, es decir $O(n)$. Al igual que la memorizacion podemos usar otros dos conceptos de la programacion dinamica:
\subsection*{Push}
Para entender esta parte de programación dinámica podemos tomar en cuenta el ejemplo de fibonacci de nuevo, tomemos en cuenta el elemento $fib(i)$, si nos damos cuenta este número contribuye a la suma de los dos siguientes elementos es decir $fib(i+1)$ y $fib(i+2)$.
\begin{lstlisting}
void push_fib(){
    vector<ll>dp(N+2,0); 
    dp[0]=0;
    dp[1]=1;
    for(int i=0;i<N;i++){
        if(i+1<=N)dp[i+1]+=dp[i];
        if(i+2<=N)dp[i+2]+=dp[i];
    }
}
\end{lstlisting}
Un ejemplo donde el enfoque \textit{Push} brilla es el problema de los saltos: Si una rana está en la posición $i$ y puede saltar a $i+1, i+2$ o $i+3$, es mucho más natural "empujar" el valor actual a los tres destinos posibles que intentar calcular quiénes llegan al destino actual.
\begin{lstlisting}
vector<ll>dp(N+4,0);
dp[0]=1;
for(int i=0;i<N;i++){
    if(dp[i]==0)continue;
    dp[i+1]+=dp[i];
    dp[i+2]+=dp[i];
    dp[i+3]+=dp[i];
}
\end{lstlisting}
\subsection*{Pull}
El enfoque \textit{Pull} es el más intuitivo y común en programación dinámica. En este caso, nos paramos en el estado actual $i$ y "traemos" (pull) los resultados de subproblemas que ya han sido resueltos previamente para calcular el valor presente.
\begin{lstlisting}
void pull_fib() {
    vector<ll>dp(N+1,0);
    dp[0]=0;
    dp[1]=1;
    for(int i=2;i<=N;i++) {
        dp[i]=dp[i-1]+dp[i-2];
    }
}
\end{lstlisting}
\end{document}